\section{Weakly asymptotically flat data in three dimensions}
\label{sec:weak3}\label{n3:intro:section}\label{n3:reduction:section}

The prepared theorem can be applied under only two differentiated metric
bounds and one differentiated tensor bound in dimension three. The reason
is a separate end construction: first match the original ADM vector by a
Schwarzschild reference slice, then bend that slice to rest and repair the
small constraint error. The metric changes at distances tending to infinity.
We prove compactness of minimizing enclosures before making those changes;
it gives convergence of the full enclosing infimum, not only a one-sided
area bound. A positive conformal family subsequently removes timelikeness
as an assumption, and truncation of the other ends gives the complete-data
statement.

\subsection{Initial data, normals, and enclosing area}

We work in three spatial dimensions, with zero cosmological constant and
\(G=c=1\). In this section $\mu,J$ denote the physical densities:
the geometric densities of Equation~\eqref{eq:setting:constraints} are
$8\pi\mu,8\pi J$. The ADM energy and momentum retain exactly their
previous normalization. We spell out this dictionary again when applying
the common prepared theorem. For a Riemannian metric \(g\) and a symmetric covariant
two-tensor \(K\), define the constraint densities by
\begin{equation}\label{n3:intro:constraints}
 16\pi\mu=R_g+(\tr_gK)^2-|K|_g^2,\qquad
 8\pi J=\Div_g\bigl(K-(\tr_gK)g\bigr).
\end{equation}
Here \(J\) is a covector field. The dominant energy condition is
\begin{equation}\label{n3:intro:dec}
                         \mu\ge |J|_g.
\end{equation}
All initial data considered here are smooth, including at a compact
boundary when one is present.

On an asymptotically flat end with Euclidean coordinates \(x\) and
\(r=|x|\), our basic assumptions are
\begin{equation}\label{n3:intro:decay}
 g_{ij}-\delta_{ij}=O_2(r^{-q}),\qquad
 K_{ij}=O_1(r^{-1-q}),\qquad q>\tfrac12.
\end{equation}
The notation \(O_j(r^{-a})\) includes the corresponding coordinate
derivative bounds through order \(j\). We assume that \(\mu\) and
\(|J|_g\) are integrable and that the following ADM limits exist and are
finite:
\begin{align}
 E&=\frac1{16\pi}\lim_{r\to\infty}
   \int_{S_r}(\partial_jg_{ij}-\partial_i g_{jj})n^i
                    \dd A_\delta,\label{n3:intro:energy}\\
 P_i&=\frac1{8\pi}\lim_{r\to\infty}
   \int_{S_r}\bigl(K_{ij}-(\tr_gK)g_{ij}\bigr)n^j
                    \dd A_\delta.\label{n3:intro:momentum}
\end{align}
The normal and area form in these definitions are Euclidean. Whenever
\(E>|P|_\delta\), write
\begin{equation}\label{n3:intro:rest-mass}
                         m=\sqrt{E^2-|P|_\delta^2}.
\end{equation}

For a two-sided surface \(\Sigma\) with specified unit normal \(\nu\),
we use
\[
 H_\Sigma=\Div_\Sigma\nu,\qquad
 \theta_+(\Sigma)=H_\Sigma+\tr_\Sigma K.
\]
Thus Euclidean spheres have positive mean curvature for the normal
toward infinity. A future marginally outer trapped surface, abbreviated
MOTS, satisfies \(\theta_+=0\). A weakly future outer trapped surface
satisfies \(\theta_+\le0\). Our spacetime sign convention is
\begin{equation}\label{n3:intro:k-sign}
 K(Y,Z)=\mathbf g(\mathbf\nabla_Y n,Z),
\end{equation}
where \(n\) is the future unit timelike normal. Equivalently, in Gaussian
normal coordinates the spatial metric has normal derivative \(2K\).

\begin{definition}[Exterior and enclosing area]\label{n3:intro:exterior-class}
An exterior is a connected orientable smooth three-manifold \(\Omega\)
with nonempty compact smooth boundary, complete as a metric space with
its boundary included, and with exactly one end, which is asymptotically flat.
An enclosing cut is the entire compact smooth embedded two-sided intrinsic
boundary of a connected closed outer domain containing the sufficiently
distant end, with its manifold interior in the open exterior. Equivalently,
it separates the entire inner boundary from infinity, with bounded
complementary pockets filled on the inner side. Every frontier component is
counted, including portions coinciding with the inner boundary. A cut may
have several components. Its normal points toward the retained end. We put
\[
 a_g(\partial\Omega)=
  \inf\{|\Gamma|_g:\Gamma\text{ is an enclosing cut}\}.
\]
When using perimeter compactness, we take the closure of this class
among filled inner sets and include the area of any frontier coinciding
with the obstacle. Smooth outward approximation gives the same
infimum. Enclosing cuts and minimizing sets are taken with bounded
complementary pockets filled. Auxiliary perimeter arguments may use
competitors containing the inner obstacle before filling: filling a bounded
pocket deletes its frontier, cannot increase perimeter, and leaves the enclosing
infimum unchanged.
\end{definition}

The definition permits coincidence because first variations of the
minimum need not be realized by a cut lying strictly outside the
obstacle. It permits disconnected cuts because neither the minimizing
hull nor an intermediate trapped boundary need be connected. The
perimeter formulation and its compactness properties are established
in Lemma~\ref{n3:reduction:area-compactness}.

\subsection{The one-ended exterior inequality}

\begin{theorem}[Minimum-enclosing-area inequality]\label{n3:intro:exterior-inequality}
Let \((\Omega,g,K)\) be an exterior as in
Definition~\ref{n3:intro:exterior-class}. Assume
Equations~\eqref{n3:intro:constraints}--\eqref{n3:intro:decay}, integrability
of \(\mu\) and \(|J|_g\), and the finite ADM limits
\eqref{n3:intro:energy}--\eqref{n3:intro:momentum}. Orient \(\partial\Omega\)
by the normal into \(\Omega\), and suppose
\(\theta_+(\partial\Omega)\le0\). If \(E>|P|_\delta\), then
\begin{equation}\label{n3:intro:penrose}
          \sqrt{E^2-|P|_\delta^2}
               \ge \sqrt{\frac{a_g(\partial\Omega)}{16\pi}}.
\end{equation}
No extension of the data across \(\partial\Omega\) is required.
\end{theorem}


\subsection{A positive lower bound for every enclosing cut}

Before preparing the end, we record a geometric fact about the full
enclosing infimum. A thin tube joining the inner boundary to infinity
carries a fixed flux through every enclosing cut. A bound for the
flux density then gives a uniform positive area bound. This argument
uses neither the constraint equations nor a sign for the boundary
expansion.

\begin{lemma}[Positive enclosing area]
\label{bridge:bridge:positive-area}
Let $(\Omega,g)$ be a connected orientable smooth three-dimensional exterior,
complete with its nonempty compact smooth boundary $B$ included, with
one asymptotically flat end and compact complement. Let $a_g(B)$ be
the infimum of the full areas of the enclosing cuts in
Definition~\ref{n3:intro:exterior-class}, including all components and
all portions coincident with $B$. Then
\[
                         0<a_g(B)<\infty.
\]
\end{lemma}

\begin{proof}
A sufficiently large coordinate sphere encloses $B$, so the infimum
is finite. To obtain a lower bound, choose a point of $B$ and a smooth
proper embedded ray starting there, initially normal to $B$ and pointing
into $\Omega$. Choose the ray to be straight in the asymptotic coordinates
outside a compact set. A tubular neighborhood has product coordinates
$(t,y)\in[0,\infty)\times D^2$, with its initial disk contained in $B$.
On the straight tail the transverse coordinate disk can have a fixed
positive width.

Choose a smooth two-form $\eta$ supported in the interior of $D^2$ with
$\int_{D^2}\eta=1$. Pull it back by $(t,y)\mapsto y$ and extend it by
zero across the lateral boundary of the tube. This gives a smooth
closed two-form $\omega$ on $\Omega$, smooth also at $B$. Its pointwise
comass is bounded: this follows from compactness on the initial portion
of the tube and from asymptotic flatness and fixed coordinate width on
the tail. Write
\[
 C_\omega=\sup_{x\in \Omega}\sup_{
       v,w\ \text{orthonormal in }(T_x\Omega,g)}|\omega_x(v,w)|<\infty.
\]
Orient the transverse disk so that the flux through a sufficiently
large outward-oriented coordinate sphere $S_R$ is one.

Let $\Gamma$ be a smooth enclosing cut lying in the interior of $\Omega$,
oriented toward the end, and choose $S_R$ outside it. The region between
$\Gamma$ and $S_R$ is compact and disjoint from $B$; all of its inner
boundary components are counted in $\Gamma$. Stokes' theorem gives
\[
       \int_\Gamma\omega=\int_{S_R}\omega=1,
       \qquad 1\le C_\omega\operatorname{Area}_g(\Gamma).
\]
This calculation also applies when the cut is disconnected. If a
smooth cut touches or coincides with $B$, push it into the interior
by a smooth collar flow pointing from $B$ into $\Omega$, filling the collar on the obstacle
side. The resulting surfaces are enclosing cuts and their areas
converge to the full area of the original cut. Thus the same lower
bound holds for boundary-coincident cuts. Taking the infimum yields
$a_g(B)\ge C_\omega^{-1}>0$. The smooth/perimeter equivalence in the
definition of the full enclosing infimum transfers this bound to its
filled-perimeter formulation as well.
\end{proof}

Throughout the construction, write $B=\partial\Omega$.
Thus $\Omega$ is smooth, connected and orientable, complete including its
nonempty compact smooth boundary, and has a single Euclidean coordinate
end with compact complement. The boundary normal points into $\Omega$.
The differentiated asymptotic assumptions are
\begin{equation}\label{n3:reduction:original-decay}
 g-\delta=O_2(r^{-q}),\qquad K=O_1(r^{-1-q}),\qquad q>\tfrac12.
\end{equation}
The density, finite-charge and weak-trapping hypotheses are those stated
in Theorem~\ref{n3:intro:exterior-inequality}. The preparation assumes
$E>|P|$; the conformal construction immediately below does not.

\subsection{A strict conformal direction}

Conformal strictification by a spacetime Poisson equation was developed
by Jaracz~\cite{Jaracz2025StrictDEC}. Here a regularized drift equation
with nonzero inner Neumann data also gives strict trapping, an exact
energy slope, unchanged momentum, and comparison of all enclosing cuts.
These quantitative conclusions are established below.

\begin{lemma}[Conformal constraint formulas]
\label{n3:reduction:conformal-formulas}
For a smooth function $f$, set
\[
 g_f=e^{4f}g,\qquad K_f=e^{2f}K.
\]
The corresponding constraint densities satisfy
\begin{align}
 16\pi e^{4f}\mu_f
   &=16\pi\mu-8\Delta_g f-8|df|_g^2,
       \label{n3:reduction:conformal-energy}\\
 8\pi J_f
   &=e^{-2f}\bigl(8\pi J+4K(\nabla f,\cdot)\bigr).
       \label{n3:reduction:conformal-momentum}
\end{align}
Here the second identity is an identity of covectors.  Consequently,
\begin{align}
 16\pi e^{4f}(\mu_f-|J_f|_{g_f})
 \ge{}&16\pi(\mu-|J|_g)\notag\\
 &+8\bigl(-\Delta_g f-|df|_g^2-|K|_g|df|_g\bigr).
       \label{n3:reduction:conformal-dec}
\end{align}
On the fixed boundary,
\begin{equation}
 \theta_{+,f}=e^{-2f}(\theta_++4\partial_\nu f).
       \label{n3:reduction:conformal-expansion}
\end{equation}
Equivalently, if $u=e^f>0$, the last summand in
Equation~\eqref{n3:reduction:conformal-dec} is
$8u^{-1}(-\Delta_g u-|K|_g|du|_g)$.
\end{lemma}

\begin{proof}
Apply the pointwise identities in Equation~\eqref{eq:end:conformal}
with $n=3$, $k=2$, $p=1$, and $s=2f$. The densities in that identity
are $\mu_{\mathrm A}=8\pi\mu$ and $J_{\mathrm A}=8\pi J$ in this
section's physical normalization. Setting the vector argument to zero
gives $8\pi e^{4f}\mu_f=8\pi\mu-4\Delta f-4|df|^2$;
multiplication by two is
Equation~\eqref{n3:reduction:conformal-energy}. The covector formula
proved there gives
$8\pi J_f=e^{-2f}(8\pi J+4K(\nabla f,\cdot))$.
Its new-metric norm introduces a further factor $e^{-2f}$, so the
triangle inequality gives Equation~\eqref{n3:reduction:conformal-dec}
with precisely its factor $8$. The boundary formula has
$k\partial_\nu s=4\partial_\nu f$ and proves
Equation~\eqref{n3:reduction:conformal-expansion}. Finally
Equation~\eqref{eq:end:conformal-U} uses $U=e^f=u$ here and gives
the stated positive-factor form. These are pointwise identities;
no decay or additional differentiated end hypothesis is used.
\end{proof}

\begin{lemma}[Strict direction under the original decay]
\label{n3:reduction:strict-direction}
Choose
\[
 0<\delta<\beta<\min(q,1).
\]
There are smooth positive functions $w$ and $\phi$ on $\Omega$, with
$w=r^{-3-\delta}$ on the end, such that
\begin{align}
 &\partial_\nu\phi=-1\quad\hbox{on }B,
       &\phi&=O_2(r^{-1}),\label{n3:reduction:strict-boundary}\\
 &-\Delta_g\phi-|K|_g|d\phi|_g\ge w,
       &\frac{-\Delta_g\phi}{w}&\longrightarrow1.
       \label{n3:reduction:strict-source}
\end{align}
The Euclidean normal flux
\[
 L_\phi=\lim_{r\to\infty}
          \int_{S_r}\partial_r\phi\,\dd A_\delta
\]
is finite and negative.  For every finite $s\ge0$, set
\begin{equation}
 \label{n3:reduction:linear-family}
 u_s=1+s\phi,\qquad g_s=u_s^4g,\qquad K_s=u_s^2K.
\end{equation}
These data are complete up to $B$ and satisfy DEC, strictly for $s>0$.
If $\theta_+\le0$ on $B$, their boundary expansion remains nonpositive
and is strictly negative for $s>0$.  Their decay exponent can be taken to be
$\min(q,1)>1/2$, their constraint densities are integrable, and
\begin{equation}
 \label{n3:reduction:strict-charges}
 E_s=E+sA_\phi,\qquad P_s=P,
 \qquad A_\phi=-\frac{L_\phi}{2\pi}>0.
\end{equation}
Their full enclosing infima satisfy
\begin{equation}
 \label{n3:reduction:strict-areas}
 a_g(B)\le a_{g_s}(B)
       \le (1+s\|\phi\|_\infty)^4a_g(B).
\end{equation}
In particular $E_s\to E$ and $a_{g_s}(B)\to a_g(B)$ as $s\downarrow0$.
No sign assumption on the ADM vector is needed.

If $K$ is compactly supported and the metric has symbol estimates of all
orders with $g-\delta=O(r^{-1})$, the function can instead be chosen with
$\phi=O_k(r^{-1})$ for every $k$, while retaining
Equations~\eqref{n3:reduction:strict-boundary} and
\eqref{n3:reduction:strict-source}.
\end{lemma}

\begin{proof}
Choose smooth $b\ge|K|_g$, positive everywhere and equal to
$b_0r^{-1-\beta}$ sufficiently far out, where $b_0>0$ is chosen large
enough.  This is possible because $\beta<q$.  Choose a smooth positive
regularizer $\zeta$, equal to $r^{-2}$ on the end, and extend $w$ smoothly
and positively to $\Omega$.  We solve
\begin{equation}
 \label{n3:reduction:drift-poisson}
 -\Delta_g\phi=b\sqrt{|d\phi|_g^2+\zeta^2}+w,
 \qquad \partial_\nu\phi=-1,
 \qquad \phi\longrightarrow0.
\end{equation}
The regularizer makes this equation smooth even at critical points.

First truncate at a large coordinate sphere $S_R$ and prescribe
$\phi=0$ there.  The Neumann and Dirichlet conditions occur on disjoint
boundary components.  For $0\le t\le1$, replace $b$ in
Equation~\eqref{n3:reduction:drift-poisson} by $tb$.  At $t=0$, the mixed
Laplacian is invertible: its Dirichlet face gives the Poincare inequality,
and the weak solution obtained from its coercive quadratic form is smooth
by boundary regularity.  At any solution, the linearization is a
Laplacian with a smooth drift of norm at most $b$ and with the homogeneous
mixed boundary conditions.  The strong maximum principle and the Hopf
boundary principle make its kernel zero.  It is a Fredholm perturbation
of the mixed Laplacian of index zero, and is therefore invertible.
The boundary estimates used here are the ordinary elliptic estimates
for Dirichlet and normal Neumann conditions: positive definiteness of
$g$ verifies their principal complementing condition, and the disjoint
boundary components can be covered separately
\cite{AgmonDouglisNirenberg1959}.  The maximum principles and local
elliptic estimates below are used in their usual uniformly elliptic
forms \cite{GilbargTrudinger}.

To close continuation at $t=1$, first work on a fixed truncation.
The elliptic $W^{2,p}$ estimate
and gradient interpolation give
\[
 \|\phi\|_{W^{2,p}}
 \le C_p\bigl(1+\|\phi\|_{L^p}+\|d\phi\|_{L^p}\bigr)
 \le \tfrac12\|\phi\|_{W^{2,p}}
       +C'_p(1+\|\phi\|_{L^p}).
\]
The constants include the fixed boundary data and are uniform in $t$.
If the supremum norms of solutions on this truncation were unbounded,
divide them by their supremum norms.  The preceding estimate with
$p>3$, compactness, and the equation give a nonzero limit $v$ with
homogeneous mixed data satisfying
\[
 -\Delta_gv=t_*b|dv|_g.
\]
This can be written $\Delta_gv+A\cdot dv=0$ with a bounded measurable
drift: take $A=t_*b\nabla v/|dv|$ off the critical set and zero on it.
The strong and Hopf maximum principles exclude a nonzero solution with
the homogeneous mixed conditions.  Thus the supremum norms are bounded.
The $W^{2,p}$ estimate, Schauder estimates, and the smooth equation now
give the bounds needed for closedness of continuation.  Existence on
each truncation follows.

Every such solution is nonnegative.  A negative minimum cannot be
interior because the right side of
Equation~\eqref{n3:reduction:drift-poisson} is positive.  At an inner-boundary
minimum, the outward-domain derivative would be negative by the Hopf
principle, whereas the prescribed derivative in that direction is $1$.
The outer value is zero.

It remains to make these bounds independent of $R$.  For fixed sufficiently large
$A$ and then sufficiently large $r_0$, the positive radial function
\[
 v_0(r)=r^{-1}(1-Ar^{-\delta})
\]
satisfies, on $r\ge r_0$,
\begin{align*}
 -\Delta_gv_0-b|dv_0|_g
 &=A\delta(1+\delta)r^{-3-\delta}
       +O(r^{-3-q})+O(r^{-3-\beta})\ge c_0w.
\end{align*}
The constants are fixed independently of the truncation.  Since
$b\zeta=o(w)$, a sufficiently large multiple of $v_0$ is a supersolution
of the regularized equation.  The comparison principle gives
\begin{equation}
 \label{n3:reduction:core-barrier}
 0\le\phi(x)\le C(1+M_R)v_0(r),\qquad
 M_R=\sup_{\Omega\cap\{r\le r_0\}}\phi,
 \quad r_0\le r\le R.
\end{equation}
The notation for the compact core includes all points outside the end
chart.  Were $M_R$ unbounded along expanding truncations, the normalized
solutions $\phi/M_R$ would have locally uniform $W^{2,p}$ bounds, including
at $B$, by the preceding local estimates and
Equation~\eqref{n3:reduction:core-barrier}.  A subsequence would converge
locally in $C^1$ to a nonnegative function with maximum one on the core,
homogeneous inner Neumann data, and a bounded-drift homogeneous equation.
Equation~\eqref{n3:reduction:core-barrier} makes its value tend to zero at
infinity.  It consequently attains a positive maximum, contradicting the
strong or Hopf principle for that homogeneous equation.  This proves the
uniform core bound.  The contradiction uses the homogeneous bounded-drift
equation.

Local compactness and diagonal extraction give a smooth solution of
Equation~\eqref{n3:reduction:drift-poisson}, with $\phi=O(r^{-1})$.
The claimed derivative count can be checked directly on annuli.  Put
\[
 \phi_{\rho}(y)=\rho\phi(\rho y),\qquad
 g_{\rho}(y)=g(\rho y),\qquad 1<|y|<4.
\]
On smaller fixed annuli the equation is
\[
 -\Delta_{g_{\rho}}\phi_{\rho}
 =\rho b(\rho y)
       \sqrt{|d\phi_{\rho}|_{g_{\rho}}^2+
                     (\rho^2\zeta(\rho y))^2}
       +\rho^3w(\rho y).
\]
Its coefficients have uniform $C^{1,1}$ bounds from the two derivatives
in Equation~\eqref{n3:reduction:original-decay}; the scaled drift is
$O(\rho^{-\beta})$ and the scaled source is $O(\rho^{-\delta})$.
The bounded zeroth-order norm, interior $W^{2,p}$ estimates, and gradient
interpolation first give uniform $C^{1,\alpha}$ bounds for some
$\alpha>0$.  The right side is then uniformly $C^\alpha$.  Schauder
estimates give uniform $C^{2,\alpha}$ bounds on smaller annuli.  This
proves $\phi=O_2(r^{-1})$ without using a third derivative of $g$ or a
second derivative of $K$.

The right side of Equation~\eqref{n3:reduction:drift-poisson} is
$w+O(r^{-3-\beta})$.  It is integrable and proves
Equation~\eqref{n3:reduction:strict-source}.  The divergence theorem, with
outward-domain normal $-\nu$ on $B$, gives
\begin{equation}
 \label{n3:reduction:phi-flux}
 \lim_{r\to\infty}\int_{S_r}\partial_{n_g}\phi\,\dd A_g
 =-|B|_g-\int_{\Omega}
       \bigl(b\sqrt{|d\phi|_g^2+\zeta^2}+w\bigr)\,\dd V_g.
\end{equation}
The Euclidean flux differs by $O(r^{-q})$ and has the same finite limit.

For every fixed finite $s\ge0$, the factor $u_s$ in
Equation~\eqref{n3:reduction:linear-family} is bounded and at least one,
so the transformed data are complete up to $B$.
The $u$-form of Lemma~\ref{n3:reduction:conformal-formulas} gives
\begin{align*}
 16\pi u_s^4(\mu_s-|J_s|_{g_s})
 &\ge16\pi(\mu-|J|_g)
        +\frac{8s}{u_s}(-\Delta_g\phi-|K|_g|d\phi|_g)
 \ge \frac{8s}{u_s}w,\\
 \theta_{+,s}
 &=u_s^{-2}\left(\theta_+-\frac{4s}{u_s}\right).
\end{align*}
This proves the asserted signs for all finite $s$, without a smallness
restriction.

Since $u_s-1=O_2(r^{-1})$, the transformed decay exponent is
$\min(q,1)>1/2$.  The exact energy transformation is
\[
 16\pi u_s^4\mu_s=16\pi\mu-8s u_s^{-1}\Delta_g\phi.
\]
Here $\Delta_g\phi=O(r^{-3-\delta})$, and the additional current term
is bounded by a fixed-$s$ constant times
$|K|_g|d\phi|_g=O(r^{-3-q})$.  Both are integrable in dimension three.
The bounded positive conformal factor also preserves integrability in
the transformed volume measure.

The leading metric change is $4s\phi\delta$, whose ADM energy flux
is $-8s\int_{S_r}\partial_r\phi\,\dd A_\delta$ before division by
$16\pi$.  All remaining metric flux terms are
$O_s(r^{-q})+O_s(r^{-1})$.  Writing $\pi=K-(\tr_gK)g$, one has
$\pi_s=u_s^2\pi$, so the change in its Euclidean momentum flux is
$O_s(r^2r^{-1}r^{-1-q})=O_s(r^{-q})$.
This proves Equation~\eqref{n3:reduction:strict-charges};
Equation~\eqref{n3:reduction:phi-flux} gives $A_\phi>0$.
No $1/r$ coefficient expansion for $\phi$ is needed.
Finally, the cut class is unchanged and every cut has area
$\int_\Gamma u_s^4\,\dd A_g$.  Bounding $u_s$ between $1$ and
$1+s\|\phi\|_\infty$ and taking infima proves
Equation~\eqref{n3:reduction:strict-areas}, including coincident and
disconnected cuts.

For the final assertion, take $b$ nonnegative, compactly supported, and
at least $|K|$; the proof is unchanged.  The end equation is then simply
$-\Delta_g\phi=w$.  Once the two-derivative estimate has been established,
successive annular elliptic estimates use the assumed symbol bounds for
$g$ and $w$ to give $\phi=O_k(r^{-1})$ for every $k$.
\end{proof}

\subsection{Compactness of enclosing minimizers}

The end replacement will change the metric on receding annuli.  The next
lemma keeps area-minimizing enclosures in a fixed compact set, where local
metric convergence controls their areas.

\begin{lemma}[Enclosing areas under receding changes]
\label{n3:reduction:area-compactness}
Let $g_j$ be smooth metrics on the fixed exterior $\Omega$, converging to $g$
locally uniformly up to $B$.  Suppose that in a fixed chart outside a
compact set they satisfy, with constants independent of $j$,
\begin{equation}
 \label{n3:reduction:area-geometry}
 c\delta\le g_j\le C\delta,\qquad r|\partial g_j|_\delta\le C.
\end{equation}
Suppose also that each $g_j$ has an asymptotically flat end, possibly
in a different chart, and admits a fixed compact enclosing competitor
with uniformly bounded area.  Then minimizing filled enclosures exist,
their boundaries lie in one fixed compact subset of $\Omega$, and
\[
 a_{g_j}(B)\longrightarrow a_g(B).
\]
The perimeter and smooth enclosing infima agree.  Each minimizing
boundary is $C^{1,1}$, and is smooth minimal away from contact with $B$.
\end{lemma}

\begin{proof}
For purposes of perimeter minimization only, extend the metrics through
a collar of $B$ and attach a fixed compact filling.  The filling need not
satisfy any constraint or curvature condition.  Extend $g_j$ so that the
extensions converge uniformly on this fixed collar and filling.  A filled
enclosure is then a bounded set containing the prescribed inner obstacle.

For each fixed metric, sufficiently large spheres in its asymptotically
flat chart have strictly positive outward mean curvature.  Clipping a
bounded enclosure at any such sphere cannot increase perimeter.  To see
this, integrate the divergence of the outward unit normal field of the
sphere foliation over the part of the enclosure outside the sphere.  Its
divergence is positive.  Its flux through the exterior boundary is at
most the area of that boundary, whereas its flux on the cutting sphere
is the negative of the new cutting area.  Thus the discarded boundary
area is at least the inserted area.  Approximation gives the same
statement for finite-perimeter sets.  The direct method on the resulting
bounded domain now gives a minimizer.  The clipping observation makes
this a minimizer against every bounded competitor, not merely those
inside the chosen truncation.

Off the obstacle the boundary is locally perimeter minimizing.  If it
contains a point of radius $r$ sufficiently large, the coordinate ball
of radius $c_1r$ about that point avoids the obstacle.  After scaling
this ball to unit size, Equation~\eqref{n3:reduction:area-geometry} gives
uniform ellipticity and uniform first-derivative bounds for its area
integrand.  The local perimeter density estimate gives
\[
 \operatorname{Area}_{g_j}
       (\partial F_j\cap B_{c_1r})\ge c_2r^2.
\]
Here $c_1,c_2>0$ are independent of $j$ and $r$. To justify the
uniform lower bound, work in a rescaled coordinate ball avoiding the
obstacle, and let $v(s)=|F_j\cap B_s|_\delta$. Removing $F_j\cap B_s$
is an admissible local comparison. For almost every $s$, minimality
and uniform metric comparability give
\[
 P_\delta(F_j;B_s)\le C v'(s),\qquad
 P_\delta(F_j\cap B_s)\le C'v'(s).
\]
The Euclidean isoperimetric inequality therefore yields
$v(s)^{2/3}\le C''v'(s)$. At a point in the support of the minimizing
frontier, both phases meet every ball. Integrating this differential
inequality gives $v(s)\ge c s^3$. Filling the complement in $B_s$
gives the same estimate for $|B_s\setminus F_j|_\delta$. The relative
isoperimetric inequality now gives
$P_\delta(F_j;B_s)\ge c's^2$ on a smaller fixed ball. Returning to
the original scale proves the displayed area bound. The argument also
applies at all support points by taking limits of reduced-boundary
points. No uniform curvature or higher metric derivative bound is used.
The fixed competitor bounds the total area independently of $j$.
Consequently the minimizing boundaries lie in a fixed compact set.

The local obstacle regularity theorem for pure area minimization with
a $C^2$ obstacle gives $C^{1,1}$ regularity, and smoothness off contact,
in dimension three; its local estimates require only the appropriate
bounded geometry of the smooth metric
\cite[Regularity Theorem 1.3]{HuiskenIlmanen2001}.
This applies to the fixed smooth obstacle and each of the smooth
metrics $g_j$ and $g$ after the collar extension. Push a minimizing
boundary slightly away from its contact set by a smooth vector field agreeing
there with the outward obstacle normal.  In a small obstacle collar its
flow increases the signed distance to the obstacle.  The resulting
$C^1$ embedded boundary is a positive distance from the obstacle.
Smooth local graph approximation, within that distance, preserves
enclosure and changes area by an arbitrarily small amount.  This proves
equality of the smooth and perimeter infima for each metric.

Enlarge the fixed compact region so that its complement is a connected
coordinate tail. Each minimizing inner set is either empty or the whole
tail there, since its frontier is absent. Boundedness excludes the latter,
so the sets themselves also lie in a common compact set.
Choose a fixed smooth auxiliary metric on this compact set. Uniform
metric comparability gives uniformly bounded auxiliary perimeters, and
BV compactness gives a limiting filled enclosure. The obstacle is still
contained in the limit. Fill any newly bounded complementary pockets;
this only removes frontier. Uniform metric convergence makes the area
integrands converge uniformly relative to the auxiliary metric, so lower
semicontinuity implies
\[
 a_g(B)\le\liminf_j a_{g_j}(B).
\]
Conversely, using a fixed enclosure with $g$-area arbitrarily close to
$a_g(B)$ gives the reverse upper-limit inequality.
\end{proof}

\subsection{Schwarzschild reference charges}

The area infimum is now stable under uniformly controlled changes on
receding annuli. We next arrange those changes. The vacuum reference end
must first have the original ADM charges; only after matching can we bend
it to a rest slice without paying a nonvanishing constraint-repair cost.

\begin{lemma}[A reference end with prescribed timelike charges]
\label{n3:reduction:boosted-reference}
Let $E>|P|$ and $m=(E^2-|P|^2)^{1/2}$.  There is a spacelike plane
near spatial infinity of Schwarzschild spacetime of mass $m$ whose
induced data, in asymptotically Euclidean coordinates on the plane,
have ADM charges $(E,P)$.  The reference data satisfy
\[
 g_*-\delta=O_k(r^{-1}),\qquad K_*=O_k(r^{-2})
\]
for every derivative order $k$.
\end{lemma}

\begin{proof}
Apply Lemma~\ref{lem:end:boost} with $n=3$, $k=2$, $p=1$, and
$M=m$. In static isotropic coordinates $(T,y)$ its spacetime metric is
\begin{equation}
 \label{n3:reduction:schwarzschild-spacetime}
 -\left(\frac{1-m/(2r)}{1+m/(2r)}\right)^2\dd T^2
       +\left(1+\frac{m}{2r}\right)^4\delta_{ij}\dd y^i\dd y^j,
 \qquad r=|y|.
\end{equation}
For $P\ne0$, put $v=P/E$ and rotate the boost axis into its direction;
for $P=0$, take $v=0$. Then
$\gamma=(1-|v|^2)^{-1/2}=E/m$, and the reference plane has charges
$(m\gamma,m\gamma v)=(E,P)$. The common lemma uses the same
future-normal convention for $K$, and its ADM prefactors become
\[
 \frac1{2k\omega_2}=\frac1{16\pi},\qquad
 \frac1{k\omega_2}=\frac1{8\pi}.
\]
Thus its charges have precisely the present signs and normalization;
the use of physical constraint densities here does not rescale them.

For the later bend, write the plane as $T=v\cdot y$ and set
$y=Ax$, where $A=(I-v\otimes v)^{-1/2}$. The induced Minkowski
metric in the $x$ coordinates is $\delta$. The perturbation of the
metric above from Minkowski space has symbol decay of order $r^{-1}$
on this fixed spacelike plane, and
Lemma~\ref{lem:end:boost} gives the induced metric and tensor orders
$r^{-1}$ and $r^{-2}$ through every fixed derivative order. These
estimates concern the exact reference end and impose no additional
derivative assumption on the original data.
\end{proof}

\subsection{Moments and compact linear corrections}

The reference end has the required charges, but interpolation can
violate DEC. Compact linear corrections reduce this error before the
conformal repair. This is related to localized constraint deformation
and its adjoint obstructions
\cite{CorvinoSchoen2006,ChruscielDelay2003}; those results are not being
invoked as a no-parity DEC rest-end or all-cut comparison theorem.

All operators in the next two lemmas are Euclidean.  For symmetric
covariant tensors define
\begin{equation}
 \label{n3:reduction:linear-operators}
 \mathcal Lh=\partial_i\partial_jh_{ij}-\Delta\tr h,
 \qquad
 (\mathcal Dk)_i=\partial_j(k_{ij}-(\tr k)\delta_{ij}).
\end{equation}
The formal adjoint kernel of $\mathcal L$ consists of affine functions;
the formal adjoint kernel of $\mathcal D$ consists of Euclidean Killing
fields.  Orthogonality to these kernels is the compatibility condition
for the compact inverses below.

\begin{lemma}[Compact scalar and symmetric divergence inverses]
\label{n3:reduction:compact-inverses}
Let $\mathcal A\subset\R^3$ be a fixed connected open annulus and let
$Q\Subset\mathcal A$.  Smooth sources supported in $Q$ admit the
following compact corrections, supported in a fixed larger compact
subset of $\mathcal A$.
\begin{enumerate}[label=\textup{(\roman*)}]
\item If $\int f=0$ and $\int x^if=0$ for $i=1,2,3$, there is a smooth
      symmetric tensor $h$ with $\mathcal Lh=f$ and
      $\|h\|_{W^{k+2,p}}\le C_{k,p}\|f\|_{W^{k,p}}$.
\item If $\int V\cdot Y=0$ for every Euclidean Killing field $Y$, there
      is a smooth symmetric tensor $k$ with $\mathcal Dk=V$ and
      $\|k\|_{W^{k_0+1,p}}\le C_{k_0,p}\|V\|_{W^{k_0,p}}$.
\end{enumerate}
The estimates hold for integers $k,k_0\ge0$ and $1<p<\infty$.
\end{lemma}

\begin{proof}
Apply Lemma~\ref{lem:end:compact-inverses} in dimension three, with
the prescribed source support $Q$ and the same connected annulus.
The scalar operator there is exactly $\mathcal L$ here, and its affine
moments are the four conditions in (i). The vector operator is exactly
$\mathcal D$; its translational and rotational moments are (ii).
The inverse of trace reversal in that proof is
$S\mapsto S-(\operatorname{tr}S)\delta/(n-1)$, which becomes
$S\mapsto S-(\operatorname{tr}S)\delta/2$. In particular its symmetric
correction is obtained from a skew tensor $D_{ijk}=-D_{jik}$ by
\begin{equation}\label{n3:reduction:stress-symmetrization}
 C_{ijk}=D_{ijk}-D_{ikj}-D_{jki}.
\end{equation}
The identities $C_{ijk}=-C_{ikj}$ and
$(C_{ijk}-C_{jik})/2=D_{ijk}$ cancel the skew part without changing
divergence. The two scalar inversions gain two derivatives, and the
row-divergence inversion followed by that correction gains one.
All supports and bounds are those of the fixed annulus. Thus the full
conclusions, including linear dependence and every stated Sobolev order,
are precisely the specializations of that lemma. This is an operator
statement about smooth sources; the estimates available from the original
weakly decaying data are verified separately in the annular replacement
proof below.
\end{proof}

\begin{lemma}[Integrable sources give sublinear first moments]
\label{n3:reduction:moment-estimates}
For data satisfying Equation~\eqref{n3:reduction:original-decay} and
integrable $\mu,|J|$, put $h=g-\delta$ and
$\pi=K-(\tr_\delta K)\delta$.  Then
\begin{align}
 \mathcal Lh&=16\pi\mu+O(r^{-2-2q}),\notag\\
 \partial_j\pi_{ij}&=8\pi J_i+O(r^{-2-2q}).
       \label{n3:reduction:linear-source-integrability}
\end{align}
In particular these Euclidean sources are integrable.  Their scalar
boundary fluxes paired with affine functions having zero constant term,
and their vector boundary fluxes paired with rotational Killing fields,
are $o(R)$ on coordinate spheres of radius $R$.
\end{lemma}

\begin{proof}
The scalar-curvature expansion has remainder bounded by
\[
 C\bigl(|h||\partial^2h|+|\partial h|^2+|K|^2\bigr),
\]
and the replacement of the momentum constraint by Euclidean divergence
has remainder bounded by
$C(|h||\partial K|+|\partial h||K|)$.  These are
$O(r^{-2-2q})$, using exactly the derivatives in
Equation~\eqref{n3:reduction:original-decay}.  They are integrable because
$2+2q>3$.

For any integrable function $F$ on the end and fixed $c>1$,
\begin{equation}
 \label{n3:reduction:sublinear-tail}
 \frac1R\int_{r_0<|x|<cR}|x|\,|F(x)|\,\dd x\longrightarrow0.
\end{equation}
Indeed the contribution inside any fixed radius $M$ tends to zero
after division by $R$; the remaining contribution is at most
$c\int_{|x|>M}|F|$.  Let $R\to\infty$ and then $M\to\infty$.

For an affine function $f$, the vector field
\[
 \mathcal Q_f(h)^i
  =f(\partial_jh_{ij}-\partial_i\tr h)
       -(\partial_jf)h_{ij}+(\partial_if)\tr h
\]
satisfies $\partial_i\mathcal Q_f(h)^i=f\mathcal Lh$.
For a Euclidean Killing field $Y$,
\[
 \partial_j(\pi_{ij}Y^i)=Y^i\partial_j\pi_{ij},
\]
because $\pi$ is symmetric.  Integrating from a fixed inner sphere and
using Equation~\eqref{n3:reduction:sublinear-tail} proves the asserted
$o(R)$ estimates.  This argument does not assert the existence of a
center-of-mass or angular-momentum limit and requires no parity
condition.
\end{proof}

\begin{proposition}[Annular replacement with a vanishing deficit]
\label{n3:reduction:annular-replacement}
Suppose $E>|P|$, and let $(g_*,K_*)$ be the reference end from
Lemma~\ref{n3:reduction:boosted-reference}, identified with the given end
in the common asymptotically Euclidean coordinates.  For all sufficiently
large $R$ there are smooth data $(\widetilde g_R,\widetilde K_R)$ equal to
$(g,K)$ on $r\le R$ and to $(g_*,K_*)$ on $r\ge4R$ such that
\begin{equation}
 \label{n3:reduction:annular-deficit}
 16\pi(\widetilde\mu_R-|\widetilde J_R|_{\widetilde g_R})
       \ge-\eta_RR^{-3},\qquad \eta_R\longrightarrow0.
\end{equation}
The possible negative part is supported in $R<r<4R$.
The metrics there are uniformly Euclidean comparable and have uniformly
bounded scaled first derivatives.  The construction uses no derivatives
of the original data beyond those in
Equation~\eqref{n3:reduction:original-decay} for its quantitative estimates.
\end{proposition}

\begin{proof}
Fix $\tfrac12<\beta<\min(q,1)$ and work on the annulus
$\mathcal A=\{1<|y|<4\}$ with $x=Ry$.  The scaled fields are
\[
 g_R(y)=g(Ry),\qquad k_R(y)=RK(Ry),
\]
and likewise for the reference data.  Their constraint densities are the
original densities multiplied by $R^2$.  If $h_R=g_R-\delta$, then
\begin{equation}
 \label{n3:reduction:scaled-smallness}
 \|h_R\|_{C^2(\mathcal A)}+\|k_R\|_{C^1(\mathcal A)}
 +\|h_{*,R}\|_{C^2(\mathcal A)}+\|k_{*,R}\|_{C^1(\mathcal A)}
 \le CR^{-\beta}.
\end{equation}
Choose $\chi\in C^\infty(\mathcal A)$ with $0\le\chi\le1$, equal to one near $|y|=1$ and zero
near $|y|=4$, with all its derivatives supported in a fixed compact
subannulus.  Blend both fields by this cutoff.

Write the constraints in the unnormalized form
\[
 \mathcal C(g,k)
 =\bigl(\Scal_g+(\tr_gk)^2-|k|_g^2,
             \Div_g(k-(\tr_gk)g)\bigr).
\]
Their Euclidean linear part is $(\mathcal Lh,\mathcal Dk)$.
On fields satisfying Equation~\eqref{n3:reduction:scaled-smallness}, the
nonlinear remainder is bounded in supremum norm by $CR^{-2\beta}$.
Thus the blended constraints equal the blended original constraint
sources, plus a quadratic error and the compact commutators
\begin{align*}
 f_R&=\mathcal L(\chi(h_R-h_{*,R}))
                      -\chi\mathcal L(h_R-h_{*,R}),\\
 V_R&=\mathcal D(\chi(k_R-k_{*,R}))
                      -\chi\mathcal D(k_R-k_{*,R}).
\end{align*}
They have $C^{0,1}$ norms at most $CR^{-\beta}$.  In detail the scalar
commutator applied to a symmetric tensor $h$ is
\begin{align*}
 &2\partial_i\chi(\partial_jh_{ij}-\partial_i\tr h)
       +(\partial_i\partial_j\chi)h_{ij}-(\Delta\chi)\tr h;
\end{align*}
one derivative of this expression uses at most two derivatives of $h$.
The vector commutator is $(\partial_j\chi)(k_{ij}-(\tr k)\delta_{ij})$;
one derivative uses at most one derivative of $k$.  No derivative of the
complete matter sources is needed.

Every affine moment of $f_R$ and every Killing-field moment of $V_R$
is $o(R^{-1})$.  For the constant moments, integration by parts on
$\mathcal A$ gives $R^{-1}$ times the difference between the physical
translation fluxes on its two boundary spheres, less the integrals of
the blended linear sources.  The flux difference tends to zero because
the reference and original ADM vectors agree.  Replacing the physical
momentum integrand by its Euclidean trace reversal costs
$O(R^{1-2q})$, which tends to zero.  By
Lemma~\ref{n3:reduction:moment-estimates}, the physical linear-source
integrals on this receding annulus tend to zero as well.
For a linear scalar test or a rotational vector test the corresponding
factor is $R^{-2}$.  Its physical boundary fluxes are $o(R)$ by the same
lemma, while the weighted source integrals are $o(R)$ by
Equation~\eqref{n3:reduction:sublinear-tail}.  This proves the claim.

Remove these small moments by fixed smooth bumps.  Explicitly, for a
basis $p_1,\ldots,p_4$ of affine functions and a fixed nonnegative bump
$\omega$ positive on a ball inside the annulus, the Gram matrix
\[
 G_{ab}=\int\omega p_ap_b
\]
is positive definite.  Subtract
$\omega\sum_{a,b}p_a(G^{-1})_{ab}\int p_bf_R$ from $f_R$.
For the vector moments use in exactly the same way the Gram matrix
$\int\omega Y_a\cdot Y_b$ of a basis of the six Euclidean Killing fields.
It too is positive definite, since a nonzero Killing field cannot vanish
on an open ball.  Both bump corrections have every fixed smooth norm
$o(R^{-1})$.

Apply Lemma~\ref{n3:reduction:compact-inverses} with the negative of the
moment-free commutators.  Since their $W^{1,p}$ norms are $O(R^{-\beta})$,
choosing $p>3$ gives metric and tensor corrections of size $O(R^{-\beta})$
in $W^{3,p}$ and $W^{2,p}$, respectively.  Sobolev embedding gives the
needed $C^{2,\alpha}$ and $C^{1,\alpha}$ bounds.  Their supports stay in
a fixed compact subset of $\mathcal A$.  Extend them by zero and undo
the scaling.  The metric remains positive definite for large $R$.

To check DEC, denote the original scaled sources by $(A_R,B_R)$, so
$A_R\ge2|B_R|_{g_R}$.  The reference sources vanish.  The corrected
sources therefore have the form
\[
 \bigl(\chi A_R,\chi B_R\bigr)
       +o(R^{-1})+O(R^{-2\beta})
\]
in supremum norm.  The norm of $B_R$ changes by at most
$CR^{-\beta}|B_R|\le CR^{-2\beta}$ when the metric is changed, since
the pointwise constraint formula gives $|B_R|\le CR^{-\beta}$.
All other nonlinear correction terms have the same quadratic bound.
Because $2\beta>1$, the scaled DEC deficit is $o(R^{-1})$.
Multiplication by $R^{-2}$ gives
Equation~\eqref{n3:reduction:annular-deficit}.  Outside the correction annulus
the data are either the original DEC data or the exact vacuum reference
data.  Smoothness follows from smoothness of the original fields and of
the compact inverse constructions, irrespective of the sizes of their
unneeded higher derivatives.
\end{proof}

\subsection{Bending to a rest slice and repairing the deficit}

The annular correction has reduced the possible DEC failure to
$o(R^{-3})$ on a region of radius comparable to $R$. We now work inside
the exact vacuum reference part. Bending makes the distant tensor vanish;
a radial conformal factor then repairs the earlier deficit at mass cost
$o(1)$. Both operations must retain the uniform geometry needed by
Lemma~\ref{n3:reduction:area-compactness}.

\begin{lemma}[A vacuum bend with uniform geometry]
\label{n3:reduction:vacuum-bend}
The data in Proposition~\ref{n3:reduction:annular-replacement} can be
modified farther out, inside their exact Schwarzschild part, so that
they are induced by a constant-static-time slice outside a compact set.
This modification adds no constraint deficit.  Write $(g_R,K_R)$ for
the resulting unscaled data.  In the common end
coordinates the metrics are uniformly Euclidean comparable,
and throughout all transition annuli
\begin{equation}
 \label{n3:reduction:bend-geometry}
 |\partial g_R|_\delta+|K_R|_\delta\le C/r.
\end{equation}
There is a smooth proper radius $\rho_R$ on the end which equals the original
coordinate radius on the gluing annulus and equals the rest isotropic
radius beyond the transitions, with
\begin{equation}
 \label{n3:reduction:radius-geometry}
 c r\le\rho_R\le C r,\qquad
 |d\rho_R|_{g_R}^2\ge c,
 \qquad |\Delta_{g_R}\rho_R|\le C/r.
\end{equation}
All transition annuli lie between fixed multiples of $R$; the constants
can depend on the prescribed timelike ADM vector, but not on $R$.
\end{lemma}

\begin{proof}
Use the description $T=v\cdot y$, $y=Ax$, from
Lemma~\ref{n3:reduction:boosted-reference}.  Choose $R_b=c_bR$ so large,
with $c_b$ fixed, that the region $|y|\ge R_b$ is entirely inside the
exact reference part.  Let $\psi$ be a smooth function equal to one on
$(-\infty,0]$, zero on $[1,\infty)$, and taking values in $[0,1]$.
Replace the plane by the graph
\begin{equation}
 \label{n3:reduction:bending-graph}
 T_R(y)=(v\cdot y)
           \psi\left(\frac{\log(|y|/R_b)}{L}\right).
\end{equation}
Its Euclidean gradient satisfies
\[
 |dT_R|_\delta\le |v|\bigl(1+\|\psi'\|_\infty/L\bigr).
\]
Choose a fixed $L$ making the right side strictly smaller than one.
The lapse and spatial coefficients in
Equation~\eqref{n3:reduction:schwarzschild-spacetime} converge to their
Minkowski values, so for all sufficiently large $R$ this graph is
uniformly spacelike.  In fact for positive Schwarzschild mass the lapse
is smaller than one and the spatial conformal factor exceeds one,
which only improves this estimate on the static exterior.

The graph agrees with the original plane before $R_b$ and with a
constant-time slice after $e^LR_b$.  Its first derivative is bounded,
and its second and third derivatives are $O(r^{-1})$ and $O(r^{-2})$.
The induced metric is uniformly comparable with $\delta$; differentiating
its pullback formula gives the first bound in
Equation~\eqref{n3:reduction:bend-geometry}.  The second fundamental form
uses the graph's second derivatives and ambient first derivatives,
divided by a uniformly positive spacelikeness factor, and gives the
other bound.  These estimates remain true in the $x$ coordinates because
$A$ is a fixed invertible linear map.  The constraints vanish on the
entire bent part by the Gauss--Codazzi equations in exact vacuum
Schwarzschild spacetime.

It remains to make the radius precise.  On the bent region keep
$\rho_R=|x|=|A^{-1}y|$.  Beyond the bend, where the slice is already
static, write $y=r\omega$ and put $a(\omega)=|A^{-1}\omega|$.  Choose
$R_c$ a sufficiently large fixed multiple of $R_b e^L$, and a cutoff
$\chi_0$ equal to one before zero and zero after one.  In that static
region set
\[
 \rho_R(r,\omega)
 =r\exp\left[
    \chi_0\left(\frac{\log(r/R_c)}{L_0}\right)\log a(\omega)\right].
\]
For fixed sufficiently large $L_0$ its radial derivative is bounded
below by a positive constant: differentiating in $r$ gives the positive
factor $\rho_R/r$ times
$1+L_0^{-1}\chi_0'\log a$, which is at least $1/2$.
The function and its first two Euclidean derivatives have the bounds
$\rho_R\asymp r$, $|d\rho_R|\le C$, and
$|\partial^2\rho_R|\le C/r$.  Before this interpolation, the same
bounds hold for $|A^{-1}y|$, and uniform metric comparability gives the
gradient lower bound.  Combining these facts with
Equation~\eqref{n3:reduction:bend-geometry} proves
Equation~\eqref{n3:reduction:radius-geometry}.  After the interpolation,
$\rho_R=r$ is precisely the static isotropic radius.  All radius ratios
used in the construction are fixed, independently of $R$.
\end{proof}

\begin{lemma}[A radial DEC repair of vanishing mass cost]
\label{n3:reduction:radial-repair}
Consider the data obtained from
Proposition~\ref{n3:reduction:annular-replacement} and
Lemma~\ref{n3:reduction:vacuum-bend}.  There is a smooth positive conformal
factor $u_R$, constant on the unchanged inner region and tending to one
at infinity, such that $(u_R^4g_R,u_R^2K_R)$ satisfy DEC everywhere.
Moreover
\begin{equation}
 \label{n3:reduction:repair-size}
 \|u_R-1\|_\infty=o(R^{-1}),\qquad
 |du_R|=o(R^{-2})\quad\hbox{on the transition region}.
\end{equation}
The final end is exactly Schwarzschild with $K=0$ and mass
$m+o(1)$.  The boundary expansion retains its original weak sign.
\end{lemma}

\begin{proof}
All geometry estimates below are uniform in $R$.  Enlarge the fixed
transition region slightly so that, in terms of $t=\rho_R/R$, it lies
inside an interval $(a,b)$ with $0<a<b<\infty$.  Arrange that $t\ge b$
is in the static spherical region, and that the support of the possible
deficit lies in a fixed smaller interval $I\Subset(a,b)$.  Increasing
$b$ if needed makes these assertions follow from
Lemma~\ref{n3:reduction:vacuum-bend}.
We take $a=1$: the compact supports in the annular correction leave
a fixed neighborhood of $t=1$ free of deficit.  The factor constructed
below will consequently be constant throughout the original region
$r\le R$, and is extended by that constant over the compact core of $\Omega$.

Let $\eta_R\to0$ be as in
Equation~\eqref{n3:reduction:annular-deficit}, replacing it by a positive
majorant if necessary.  We will choose $\varepsilon_R=C_0\eta_R$ and
construct
\[
 u_R=1+\frac{\varepsilon_R}{R}F_R(\rho_R/R),
\]
where $F_R$ and the relevant derivatives are bounded independently of
$R$.  Write $z_R=-F_R'$.  On $[a,b+1]$, initially define $z_R$ by
\begin{equation}
 \label{n3:reduction:repair-ode}
 z_R'=C_1z_R+\omega_1,
 \qquad z_R(a)=0,
\end{equation}
where $\omega_1$ is smooth, nonnegative, zero near $a$, and at least one
on $I$.  Take $C_1$ sufficiently large, depending only on the constants
in Equations~\eqref{n3:reduction:bend-geometry} and
\eqref{n3:reduction:radius-geometry}.  The solution is nonnegative and
has uniform bounds on this fixed interval.  For a function with
$F_R'=-z_R$ the chain rule gives
\begin{align*}
 -\Delta_{g_R}u_R-|K_R|_{g_R}|du_R|_{g_R}
 &=\frac{\varepsilon_R}{R^3}
   \left[z_R'|d\rho_R|_{g_R}^2
       +Rz_R\Delta_{g_R}\rho_R
       -Rz_R|K_R|_{g_R}|d\rho_R|_{g_R}\right]\\
 &\ge\frac{\varepsilon_R}{R^3}
       (c_1z_R'-C_2z_R).
\end{align*}
Choose $C_1$ so that this is nonnegative and is at least
$c_1\varepsilon_RR^{-3}$ on $I$.

The preceding differential inequality need not be imposed on the
spherical tail.  There $K_R=0$ and the Schwarzschild radial Laplace flux
coefficient is $(r+m/2)^2$.  Put $c_R=m/(2R)$ and, on the still spherical
interval near $b$, define
\[
 A_R(t)=(t+c_R)^2z_R(t).
\]
It has positive derivative wherever $z_R>0$.  Continue it to a positive
constant on $[b+1,\infty)$ while keeping $A_R'\ge0$.  For an explicit
smooth continuation, continue the solution of
Equation~\eqref{n3:reduction:repair-ode} to $b+1$, multiply the derivative
of $(t+c_R)^2z_R(t)$ by a nonnegative cutoff equal to one near $b$ and
zero near $b+1$, and integrate with the initial value $A_R(b)$.
Define $z_R=A_R(t)/(t+c_R)^2$ on this tail, and set it to zero for
$t\le a$.  All joins are smooth, the functions are nonnegative, and
their bounds on finite intervals are uniform in $R$.  The constant
tail value, denoted $A_R^\infty$, is bounded independently of $R$.

Now define
\[
 F_R(t)=\int_t^\infty z_R(s)\,\dd s.
\]
It is nonnegative, uniformly bounded, and constant before $a$.
The preceding construction gives a superharmonic $u_R$ on the spherical
tail, because its signed radial flux is
$-\varepsilon_RA_R(t)$ and is nonincreasing.  It also gives
Equation~\eqref{n3:reduction:repair-size}.  For large $R$, $1\le u_R\le2$.
Increasing $C_0$ ensures
\[
 -\Delta_{g_R}u_R-|K_R||du_R|
 \ge \frac{u_R}{8}\,
       \bigl[16\pi(|J_R|_{g_R}-\mu_R)\bigr]_+.
\]
On the original inner region the conformal factor is constant and the
original data satisfy DEC.  Elsewhere this inequality and
Lemma~\ref{n3:reduction:conformal-formulas} prove DEC for the repaired data.
Using $u_R$, rather than estimating $\log u_R$ separately, incorporates
the gradient-square term exactly.

On the final end,
\[
 u_R=1+\frac{\varepsilon_RA_R^\infty}{r+m/2}.
\]
Writing $a_R=\varepsilon_RA_R^\infty=o(1)$, one obtains the exact identity
\begin{equation}
 \label{n3:reduction:repaired-schwarzschild}
 u_R^4\left(1+\frac{m}{2r}\right)^4\delta
   =\left(1+\frac{m/2+a_R}{r}\right)^4\delta.
\end{equation}
Thus the final mass is $m+2a_R=m+o(1)$ and the momentum is zero.
Multiplication of $K_R$ by $u_R^2$ preserves its compact support.
On the inner boundary, where $u_R$ is constant, the expansion is simply
multiplied by $u_R^{-2}$, so its weak sign is preserved.
\end{proof}

\subsection{The numerical reduction}

The preceding construction has replaced the end while preserving DEC
and the boundary expansion sign.  A final strict conformal change puts
the data in the class used for the elliptic deformation.

\begin{proposition}[Reduction to strict data with a controlled rest end]
\label{n3:reduction:rest-reduction}
Let $(\Omega,g,K)$ satisfy the exterior hypotheses of
Theorem~\ref{n3:intro:exterior-inequality}, including $\theta_+\le0$ on its
smooth compact boundary, and let $m=(E^2-|P|^2)^{1/2}>0$.
There is a sequence of smooth exterior data $(g_j,K_j)$ on $\Omega$ with the
following properties:
\begin{enumerate}[label=\textup{(\roman*)}]
\item $\mu_j>|J_j|_{g_j}$ everywhere and $\theta_{+,j}<0$ on $B$.
      The constraint densities are integrable.  On the end their DEC
      margin is bounded below by $c_jr^{-3-\delta}$ for some $c_j>0$
      and one fixed $0<\delta<1$. For a smooth positive radius extension
      $\rho$, the constant can be decreased so that
      $\mu_j-|J_j|_{g_j}\ge c_j\rho^{-3-\delta}$ on all of $\Omega$.
\item $K_j$ is compactly supported, and in a rest chart
      $g_j-\delta=O_k(r^{-1})$ and $R_{g_j}=O_k(r^{-3-\delta})$
      for every $k$.  In particular
      $\Ric_{g_j}=O(r^{-3})$ there.  Their ADM momentum is zero.
\item $(g_j,K_j)\to(g,K)$ smoothly on every fixed compact subset up to
      the boundary, their ADM energies satisfy $E_j\to m$, and
      $a_{g_j}(B)\to a_g(B)$.
\end{enumerate}
Consequently it suffices to prove the energy/enclosing-area inequality
for the class in \textup{(i)}--\textup{(ii)} in order to prove
Theorem~\ref{n3:intro:exterior-inequality} in its stated class.
\end{proposition}

\begin{proof}
Apply Proposition~\ref{n3:reduction:annular-replacement},
Lemma~\ref{n3:reduction:vacuum-bend}, and
Lemma~\ref{n3:reduction:radial-repair} along any sequence $R_j\to\infty$.
The resulting DEC exteriors have compactly supported second fundamental
form and an exact rest Schwarzschild end with energy tending to $m$.
They agree with the original data on expanding compact subsets except
for a constant conformal factor tending to one.  Hence they converge
smoothly on each fixed compact subset.  Their boundary expansion is
weakly negative.

Apply the last part of Lemma~\ref{n3:reduction:strict-direction} separately
to each of these fixed exteriors, using the same chosen
$0<\delta<1$ and a compactly supported drift coefficient.  Choose its
positive parameter $s$ in Equation~\eqref{n3:reduction:linear-family}
small enough that the energy change, $\|\log(1+s\phi)\|_\infty$, and
the supremum of its first derivative multiplied by the fixed original
coordinate radius are at most $1/j$.  Also require the change in the
first $j$ smooth norms on the first $j$ members of a compact exhaustion
to be at most $1/j$.
All these quantities are finite for the fixed $j$-th exterior, so these
choices are possible.  No estimate uniform in the receding bending
radius is required for this final strictification.

The resulting data have pointwise strict DEC and strict trapping.
The quantitative end gap is at least $c_j r^{-3-\delta}$. On its compact
complement the continuous positive function
$(\mu_j-|J_j|_{g_j})\rho^{3+\delta}$ has a positive minimum.
Replacing $c_j$ by the smaller of this minimum and the end constant
proves the global weighted assertion in \textup{(i)}. The strict conformal factor has
symbol decay of all orders and does not enlarge the support of $K$.
On the far Schwarzschild end the drift vanishes, so
$-\Delta\phi=w=r^{-3-\delta}$. Since the rest Schwarzschild metric
has zero scalar curvature and $K=0$ there, the exact $u$-formula gives
\[
 R_{g_j}=8s(1+s\phi)^{-5}r^{-3-\delta}.
\]
This proves the scalar-curvature bound in \textup{(ii)}.
The Ricci bound follows from
the two-derivative metric decay. The energy change tends to zero and the
momentum change is zero. The prescribed smallness gives smooth local
convergence, proving \textup{(ii)} and the local convergence and charge
assertions in \textup{(iii)}.

The metrics before strictification satisfy the uniform comparisons and
scaled first-derivative bounds of
Lemma~\ref{n3:reduction:area-compactness}, in the fixed original end chart:
the reference plane, the bend, and the radius interpolation involve only
fixed linear maps and fixed radius ratios.  The radial repair preserves
these bounds.  The final strictification parameters were chosen to
preserve them as well.  That lemma therefore gives the convergence of
enclosing infima.

This completes the preparation; no numerical inequality was used in
constructing these data.
\end{proof}

\subsection{Applying the prepared theorem at each fixed stage}
\label{n3:reduction:prepared-transfer}

\begin{proof}[Proof of Theorem~\ref{n3:intro:exterior-inequality}]
Let $m=\sqrt{E^2-|P|^2}>0$, and take the sequence in
Proposition~\ref{n3:reduction:rest-reduction}. Each datum is a smooth
complete one-ended exterior on $\Omega$ with the same nonempty compact
boundary $B$. The proposition gives a compactly supported tensor,
$O_d(r^{-1})$ metric decay for every fixed derivative order $d$,
scalar curvature $O(r^{-3-\delta})$ with $0<\delta<1$, integrable
constraints, strict future trapping, and a global weighted DEC margin.
These are the geometric and asymptotic conditions in
Definition~\ref{def:prepared}.

The enclosing class also agrees with its full-cut class: both retain
a connected end-containing outer domain and count its entire intrinsic
frontier, including coincidence with the obstacle. Filling pockets
removes only surplus frontier and leaves either infimum unchanged.
Thus the prepared area is $A_j=a_{g_j}(B)>0$, where positivity follows
from Lemma~\ref{bridge:bridge:positive-area}.

There is one normalization conversion. The densities used in this section
are physical densities, while those of the prepared theorem satisfy
$2\mu_{\mathrm A}=R+(\operatorname{tr}K)^2-|K|^2$ and
$J_{\mathrm A}=\operatorname{div}(K-(\operatorname{tr}K)g)$. Hence
\[
 \mu_{\mathrm A}=8\pi\mu,\qquad J_{\mathrm A}=8\pi J.
\]
The global weighted margin becomes
$\mu_{\mathrm A,j}-|J_{\mathrm A,j}|_{g_j}
\ge8\pi c_j\rho^{-3-\delta}$, and all sign and integrability
conditions are preserved. At $n=3$, $k=n-1=2$ and $\omega_2=4\pi$,
so the energy and momentum prefactors in that theorem are respectively
\[
 \frac1{2k\omega_2}=\frac1{16\pi},\qquad
 \frac1{k\omega_2}=\frac1{8\pi}.
\]
Thus neither ADM energy nor momentum nor invariant mass is rescaled.

Theorem~\ref{thm:numerical:prepared} therefore applies to each fixed
prepared datum and gives
\[
 E_j\ge\frac12\left(\frac{a_{g_j}(B)}{4\pi}\right)^{1/2}
     =\sqrt{\frac{a_{g_j}(B)}{16\pi}}.
\]
Its elliptic exhaustion and its numerical limits in $N$ and $\epsilon$
have already been completed for that datum; no estimate uniform in
$j$ is required. Now use $E_j\to m$ and the full convergence
$a_{g_j}(B)\to a_g(B)$ to obtain
\[
 \sqrt{E^2-|P|^2}\ge\sqrt{\frac{a_g(B)}{16\pi}},
\]
as asserted.
\end{proof}

\subsection{Future timelikeness from the exterior inequality}
\label{bridge:bridge:section}

The conformal direction already constructed allows us to move the ADM
energy upward while fixing the momentum and increasing every enclosing
area. If the original ADM vector were not future timelike, we could
approach the future null cone through data to which the timelike
inequality applies. The positive area bound would persist along this
family, giving a contradiction.

\begin{theorem}[The exterior inequality without a timelike assumption]
\label{bridge:bridge:timelike}\label{thm:weak3}
Let $(\Omega,g,K)$ satisfy all the hypotheses of
Theorem~\ref{n3:intro:exterior-inequality} except $E>|P|_\delta$.
Then
\[
 E>|P|_\delta,\qquad
 \sqrt{E^2-|P|_\delta^2}\ge
                  \sqrt{\frac{a_g(B)}{16\pi}},\qquad B=\partial \Omega.
\]
No extension of the data across $B$ is required.
\end{theorem}

\begin{proof}
Lemma~\ref{n3:reduction:strict-direction} supplies a positive function
$\phi$, a constant $A_\phi>0$, and, for every finite $s\ge0$, the data
\[
       g_s=(1+s\phi)^4g,\qquad K_s=(1+s\phi)^2K.
\]
These data retain completeness, integrable constraint densities, DEC,
weak future trapping of $B$, and the required differentiated decay.
The lemma makes no sign assumption on the original ADM vector. Its
charge and area identities, Equations~\eqref{n3:reduction:strict-charges}
and~\eqref{n3:reduction:strict-areas}, give
\[
       E_s=E+sA_\phi,\qquad P_s=P,\qquad
       a_{g_s}(B)\ge a_g(B)>0,
\]
where the last strict inequality is
Lemma~\ref{bridge:bridge:positive-area}.

Suppose $E\le|P|_\delta$, and set
$s_0=(|P|_\delta-E)/A_\phi$. This is a finite nonnegative number.
For every $s>s_0$ the deformed vector is future timelike, so
Theorem~\ref{n3:intro:exterior-inequality}, applied to these fixed data,
implies
\[
 \sqrt{(E+sA_\phi)^2-|P|_\delta^2}
       \ge\sqrt{\frac{a_{g_s}(B)}{16\pi}}
       \ge\sqrt{\frac{a_g(B)}{16\pi}}>0.
\]
As $s\downarrow s_0$ the left side tends to zero, while the final
lower bound is fixed and positive. This contradiction proves
$E>|P|_\delta$. The timelike exterior inequality now applies directly
to $(\Omega,g,K)$ and gives the stated bound.
\end{proof}

\subsection{Complete data with finitely many ends}

We now pass from a one-ended exterior to a complete boundaryless
initial-data manifold. The ambient enclosing class allows other ends
to remain on the chosen side. We first isolate the chosen end to use
the exterior theorem, and then compare the two area infima in the
direction required for the conclusion.

\begin{definition}[Admissible boundary and ambient enclosing area]
\label{bridge:intro:global-admissible}
Let $(M,g,K)$ be a smooth three-dimensional initial data set without
boundary, with $M$ connected, and fix one asymptotically flat end $e$.
An admissible boundary for $e$ is a
nonempty compact smooth embedded two-sided surface $S=\partial D^+$,
where $D^+$ is a smooth open region containing the entire sufficiently
distant part of $e$. The surface and the region may be disconnected,
and $D^+$ may contain other ends. For the normal $\nu$ pointing into
$D^+$, require
\[
                  H_S+\operatorname{tr}_S K\le0.
\]
Its minimum enclosing area is
\begin{equation}
\label{bridge:intro:global-area}
 A_{\min}(S)=\inf\left\{
 \operatorname{Area}_g(\partial D'^+):
 \begin{array}{l}
 D'^+\subseteq D^+\text{ is a smooth open region containing}\\
 \text{the entire sufficiently distant part of }e,\\
 \partial D'^+\text{ is compact, smooth, and embedded}
 \end{array}\right\}.
\end{equation}
The full boundary area is counted, including portions coincident
with $S$. Competitor regions $D'^+$ may be disconnected and may retain
other ends. Their boundaries need be neither trapped nor minimal.
The definition imposes no sign condition on the inward null expansion.
\end{definition}

\begin{theorem}[Spacetime Penrose inequality for complete initial data]
\label{bridge:intro:global-penrose}
Let $(M,g,K)$ be a connected orientable smooth three-dimensional initial
data set without boundary, with $g$ complete. Suppose that outside a
compact set $M$ has finitely many asymptotically flat ends. On each end,
in its asymptotic coordinates, assume
\[
 g-\delta=O_2(r^{-q}),\qquad K=O_1(r^{-1-q}),\qquad q>\tfrac12,
\]
and finite ADM energy and momentum limits with the normalizations
\eqref{n3:intro:energy}--\eqref{n3:intro:momentum}. The decay exponents
may differ between ends. Assume that the physical constraint densities
\eqref{n3:intro:constraints} obey $\mu\ge|J|_g$ and that both $\mu$ and
$|J|_g$ are integrable on $M$.

For every chosen end and every admissible boundary $S$ in
Definition~\ref{bridge:intro:global-admissible}, its ADM vector $(E,P)$
satisfies
\begin{equation}
\label{bridge:intro:global-inequality}
 E>|P|_\delta,\qquad
 \sqrt{E^2-|P|_\delta^2}\ge
                         \sqrt{\frac{A_{\min}(S)}{16\pi}}.
\end{equation}
\end{theorem}

\begin{proof}
Let $D_0$ be the connected component of $D^+$ containing the distant
chosen end. Its boundary is a union of components of $S$. It is
nonempty: otherwise $D_0$ would be both open and closed in the connected
manifold $M$, and hence equal to $M$, contradicting $S\ne\varnothing$.
The closure $X_0=\overline{D_0}$ is therefore a smooth manifold with
nonempty compact boundary. Since $S$ is compact, each sufficiently
distant end of $M$ lies either entirely in $D_0$ or entirely outside
it. Thus $X_0$ has finitely many ends, precisely those original end
tails contained in $D_0$.

For every end of $X_0$ other than the chosen one, remove a sufficiently
distant open coordinate tail. The tails can be chosen disjoint and
disjoint from $S$. Denote the remainder by $\Omega$. It is connected: a
path segment entering one of the removed tails can be replaced by
a path on the connected coordinate sphere bounding that tail.
Consequently $\Omega$ is a one-ended smooth exterior with compact boundary
\[
             B=\partial D_0\ \cup\
               \bigcup_{\text{other ends of }X_0} S_{R_i}.
\]
This boundary is nonempty because $\partial D_0$ is nonempty.

The set $\Omega$ is closed in the complete manifold $(M,g)$, and it is
complete for its intrinsic distance with $B$ included. Indeed, an
intrinsic Cauchy sequence is ambient Cauchy and has its ambient limit
in $\Omega$. At an interior limit point the two distances are locally
comparable. At a boundary limit point, a smooth boundary chart
identifies $\Omega$ locally with a half-ball and gives the same comparison,
so the sequence converges intrinsically as well. Restricting the data
preserves DEC and integrability. The chosen end is unchanged, so its
decay and ADM vector remain the original ones.

On $\partial D_0$ the normal into $\Omega$ is the normal into $D^+$, and
the future expansion is nonpositive. On an added sphere $S_{R_i}$,
the normal into $\Omega$ points toward decreasing radius in the removed
end. The differentiated asymptotic assumptions give
\[
 H_{S_{R_i}}=-\frac{2}{R_i}+O(R_i^{-1-q_i}),\qquad
 \operatorname{tr}_{S_{R_i}}K=O(R_i^{-1-q_i}).
\]
Choose each $R_i$ large enough that their sum is negative. Every
component of $B$ is then weakly future outer trapped for the normal
into $\Omega$. Theorem~\ref{bridge:bridge:timelike} applies and yields
\begin{equation}
\label{bridge:bridge:restricted-inequality}
 E>|P|_\delta,\qquad
 \sqrt{E^2-|P|_\delta^2}\ge\sqrt{\frac{a_g(B)}{16\pi}}.
\end{equation}

It remains to compare the enclosing classes. We claim
\begin{equation}
\label{bridge:bridge:global-area-comparison}
                         a_g(B)\ge A_{\min}(S).
\end{equation}
Take first a smooth enclosing cut $\Gamma$ contained in the interior
of $\Omega$. The connected component on its distant chosen-end side is
a smooth open region $D'^+\subset D_0\subset D^+$. Its boundary
consists of components of $\Gamma$, so it is compact and smooth and
has area at most $\operatorname{Area}_g(\Gamma)$. It is therefore an
admissible competitor in \eqref{bridge:intro:global-area}.
For a cut that meets or coincides with $B$, a smooth collar push into
the interior, with the obstacle side filled, gives enclosing cuts
whose areas converge to its full area. This applies also to contact
with the added coordinate spheres. Hence every full enclosing cut
of $\Omega$ has area at least $A_{\min}(S)$, proving
\eqref{bridge:bridge:global-area-comparison} after taking the infimum.
The discarded components of $D^+$ do not obstruct this comparison:
the ambient definition permits the chosen-end component alone as
a competitor. Combining the area comparison with
\eqref{bridge:bridge:restricted-inequality} proves the theorem.
\end{proof}

The auxiliary exterior excludes the other distant ends, whereas the
ambient competitor class allows them. The proof uses only the
one-sided comparison \eqref{bridge:bridge:global-area-comparison}.
Rigidity is a separate question about the original data; it requires
its own argument rather than equality transport through these
conformal changes or end truncations.
